\section{Open questions}
%==========================================================================

\paragraph{Well-posedness.} Three quantities in the formulation are not measured and are
not physiological constants: $\bar m$ (clinical convention), $\spk$ (varies with the
precipitating event), and the map from intracellular substrate to plasma biomarker. $\thc$
is therefore an operational definition rather than a physical parameter. Whether the
decision built on it is invariant to that arbitrariness is the central question.

\paragraph{Identification.} Under what conditions is a parameter identifiable when the
excitation required to identify it is the regime that a safety constraint forbids? Under
therapy the excitation is supplied by the dosing schedule itself, but the schedule must be
selected \emph{before} any response data exist: the information arrives after the moment at
which it is required.

\paragraph{Design.} The safety constraint is active only at the trough preceding each dose,
$\theta_{\mathrm{base}} + D e^{-\lambda \Delta t} > \thc$. Designing a sequence of
administration times that guarantees this for every parameter value in an uncertainty set,
while simultaneously contracting that set, is dual control with a state constraint whose
defining parameter is the unknown. Violation is not an off-specification excursion but a
hospitalisation with permanent neurological cost, so the appropriate formulation is a chance
constraint rather than a quadratic objective.

\paragraph{Reduction.} A model reduction preserving the input--output response to periodic
impulsive inputs does not obviously yield a reduced model whose parameters are identifiable
from sparse sampling. These are distinct objectives and may conflict.

\paragraph{Multi-parameter case.} The scalar analysis above is complete. With several
unknowns entering simultaneously, $\theta$, the population parameters, $\lambda$, and one
amplitude per documented illness episode, attractivity of the estimator manifold becomes a
matrix condition rather than a sign condition, and is not addressed here.
